\documentclass[11pt,a4paper]{article}
\usepackage[T1]{fontenc}
\usepackage[utf8]{inputenc}
\usepackage{lmodern}
\usepackage[margin=25mm]{geometry}
\usepackage{microtype}
\usepackage{amsmath,amssymb}
\usepackage{graphicx}
\usepackage{booktabs,tabularx,array}
\usepackage{enumitem}
\usepackage{caption}
\usepackage{xspace}
\usepackage{xcolor}
\usepackage[numbers,sort&compress]{natbib}
\usepackage[hyphens]{url}
\usepackage[colorlinks=true,linkcolor=black,citecolor=black,urlcolor=blue]{hyperref}
\usepackage{fancyhdr}
\usepackage{placeins}
\usepackage{listings}
\hypersetup{pdftitle={Missile Intercept: Modular Simulation Software and an Audit of Its Computational Evidence},pdfauthor={Anonymous Authors}}
\definecolor{ink}{HTML}{162D44}
\captionsetup{font=small,labelfont=bf}
\setlength{\parskip}{3pt}
\setlength{\parindent}{1.2em}
\setlist{nosep,leftmargin=*}
\setlength{\emergencystretch}{2em}
\clubpenalty=10000
\widowpenalty=10000
\pagestyle{fancy}
\fancyhf{}
\fancyhead[L]{\small\textit{Missile Intercept: Software Architecture and Evidence Audit}}
\fancyhead[R]{\small Anonymous preprint}
\fancyfoot[C]{\thepage}
\renewcommand{\headrulewidth}{0.3pt}
\setlength{\headheight}{14pt}
\input{evidence_numbers}
\newcommand{\code}[1]{\texttt{#1}}
\newcolumntype{Y}{>{\raggedright\arraybackslash}X}
\title{\vspace{-1.2em}\textbf{Missile Intercept: Modular Simulation Software\\and an Audit of Its Computational Evidence}}
\author{Anonymous Authors}
\date{November 13, 2026}

\begin{document}
\maketitle
\thispagestyle{plain}

\begin{abstract}
Simulation repositories often combine source code, telemetry, and polished reports without preserving the execution history that connects them. We present an artifact-based software study of \emph{Missile Intercept}, a Python framework described as a modular six-degree-of-freedom interception simulator. The study documents its software organization and examines the evidence available in a local source archive. A static inventory records \SelectedFiles selected files and \ReportAssets local assets across three historical reports. A contemporary run of the independent provenance-tooling suite yields \TestsPassed passed tests and \TestsSkipped platform-dependent skip; \ProbesMet of \ProbeCount additional synthetic integrity probes match their expected outcomes. All three reports lack matching provenance manifests and remain \code{legacy\_unverified}, despite no missing local assets being detected within the static audit scope. The findings support a narrow conclusion about current artifact organization and general file-integrity behavior. They do not establish the provenance of historical simulation outputs, reproduce a simulation campaign, or validate physical performance. The accompanying package provides the manuscript, vector figures, machine-readable audit records, and the generic tooling used in the study.
\end{abstract}

\noindent\textbf{Keywords:} scientific software; simulation architecture; artifact provenance; reproducibility; software verification

\section{Introduction}
Scientific simulation software produces several different kinds of evidence. Source files describe an implementation, test records document observations made in a particular environment, and reports communicate selected results. These objects can coexist in one archive without establishing that a particular report was generated by a particular implementation. This gap becomes especially consequential when a report contains precise numerical claims: presentation quality can make the claim easier to read without making its origin easier to verify.

Practical guidance for research computing emphasizes organized projects, version control, and scriptable analysis and manuscript workflows \citep{wilson2017}. Reproducible computational research also requires retaining the steps, program versions, parameters, and inputs that produced a result \citep{sandve2013}. These recommendations concern relationships between artifacts, rather than the mere presence of files. A source archive downloaded after a calculation cannot, by itself, recover the calculation's original settings or runtime.

This paper studies that distinction in the \emph{Missile Intercept} archive. Its documentation describes a Python flight-simulation framework with a rigid-body model, environmental and sensor modules, orchestration, telemetry, and analysis tools. The archive also contains three historical HTML reports and independent tools for environment inventory, file snapshots, and report auditing. The root documentation explicitly identifies the historical reports as unverified and states that model parameters are engineering extrapolations rather than measurements of a fielded system.

The study asks three questions. \textbf{RQ1:} How is the archived software organized, and which interfaces connect implementation and recorded artifacts? \textbf{RQ2:} What can the available report assets and manifests establish about the historical reports? \textbf{RQ3:} What behavior can be demonstrated now for the independent artifact-audit tooling? These are software and evidence questions. We do not evaluate interception effectiveness or introduce a new guidance, control, or vehicle-design method.

The paper makes three bounded contributions:
\begin{enumerate}
  \item A consolidated architecture description, with implementation descriptions kept distinct from executed verification.
  \item A reproducible static audit of the three report bundles, including explicit local-asset coverage and provenance status.
  \item A contemporary execution record for the independent tooling tests and ten controlled, nonphysical integrity probes, supplied with a claim-to-evidence map.
\end{enumerate}
The snapshot tools and simulation software already exist in the archive. Our contribution is their documentation and empirical examination in this case study, rather than a claim to have invented a new provenance standard.

\section{Evidence and Related Work}
\label{sec:related}
\citet{sandve2013} frame reproducibility around preserving how each result was produced. \citet{wilson2017} provide complementary guidance on organizing software, data, and manuscripts so that those relationships can be maintained. The present study applies these ideas to an archive in which a historical execution chain is missing; it does not claim that documenting the current archive reconstructs that chain.

The W3C PROV data model represents provenance through entities, activities, agents, and their relationships \citep{moreau2013}. In this vocabulary, a report and a source file are entities, while a simulation run and a report-generation step are activities. A file inventory records entities and their present attributes. It does not observe the historical activities that generated them. The archive's JSON manifests are lightweight application-specific snapshots; no conformance to W3C PROV serialization or semantics is claimed.

Cryptographic digests provide a useful integrity primitive. The Secure Hash Standard specifies algorithms for detecting changes to messages through their digests \citep{nist2015}. A digest comparison supports a statement about selected bytes relative to a retained reference. Without a trusted origin for that reference, it does not authenticate who created an artifact or establish that its contents are correct.

\citet{barba2021} discuss the importance of provenance and auditability for trustworthy computational evidence, including situations in which the original computation is difficult to repeat. We use four explicit terms to avoid an ambiguous claim of ``reproducibility'': \emph{static inventory} records current files; \emph{snapshot verification} compares selected files with a manifest; \emph{execution evidence} records a newly performed check; and \emph{historical reproduction} would require executing the recorded historical workflow and assessing its outputs. Only the first three are undertaken here, and execution evidence is limited to generic tooling.

\section{Software Architecture}
\label{sec:architecture}
\subsection{Archive organization and model scope}
The archive has two relevant software layers (Figure~\ref{fig:architecture}). The \code{missile\_intercept/} directory contains the simulation package. Root-level \code{tools/} and \code{tests/} contain a separate artifact-audit implementation and its tests. These layers have different dependency requirements and different evidential roles. The static inventory identifies \SimulationPythonFiles Python files under the simulation directory, including its tests and utilities.

The simulation documentation describes translational and rotational rigid-body state, quaternion attitude, changing mass properties, environmental models, navigation and sensing, phase-specific decision modules, and control interfaces. These descriptions characterize the intended software model. They are not measurements of an actual vehicle, and their presence in source code is not evidence that a historical end-to-end calculation was correct.

Table~\ref{tab:modules} summarizes the main software responsibilities. The simulator's mathematical and operational parameters are outside the present artifact study. The analysis instead follows where state is represented, where outputs are produced, and how those outputs are preserved for later inspection.

\begin{figure}[t]
  \centering
  \includegraphics[width=\linewidth]{figures/architecture.pdf}
  \caption{Organization of the archived simulation and the independent audit layer. The upper lane is described through static source and documentation inspection; the lower lane contains the tooling executed in this study. The arrows summarize artifact organization rather than an exact runtime call graph.}
  \label{fig:architecture}
\end{figure}

\begin{table}[t]
\caption{Archived software responsibilities and their relevance to the evidence study. Module presence is a static observation, not a successful execution result.}
\label{tab:modules}
\small
\begin{tabularx}{\linewidth}{@{}p{0.27\linewidth}YY@{}}
\toprule
Component & Software responsibility & Evidence implication \\
\midrule
\code{main.py}, \code{integrator.py} & Orchestration, integration, scheduling, and event interfaces & Execution history must identify the code and configuration actually used. \\
\code{state\_schema.py}, \code{config.py} & Named state layout, typed data containers, and settings & Interface meaning must accompany raw numeric arrays. \\
\code{vehicle/}, \code{sensors/} & Environmental, vehicle-state, and observation modules & Model implementation and physical validation are separate questions. \\
\code{guidance/}, \code{control/} & Phase-specific decision and command modules & Their existence does not establish effectiveness or historical use. \\
\code{telemetry.py}, \code{visualization.py} & Recorded output and presentation & A displayed figure requires an attributable input and generation step. \\
\code{batch.py}, \code{analysis.py} & Campaign and analysis utilities & Aggregate claims require case-level execution records. \\
Root \code{tools/}, \code{tests/} & Environment metadata, snapshots, and static report audit & Can be checked independently of the simulator. \\
\bottomrule
\end{tabularx}
\end{table}

\subsection{State contracts and orchestration}
The state-schema implementation associates each named field with its shape, units, frame, and owner. It computes named slices and a hash of the serialized layout. Packing and unpacking functions check field names and array shapes. At the software level, these mechanisms make a change in interface layout inspectable and reduce dependence on unexplained array offsets. A layout hash identifies the schema description; it does not verify the physical meaning of the values stored under that schema.

The orchestration source describes a sampled-data outer loop and continuous integration between updates. Discrete module outputs are held during the internal integration stages. We record this as an implementation convention observed in source comments and interfaces. We do not execute the simulation to measure numerical convergence, scheduler behavior, or the consequences of that convention.

Telemetry is collected through a preallocated buffer and written as named columns. The logger includes checks on frame width, finite values, storage capacity, and time ordering. Such guards can expose certain data-contract violations, but a well-formed output row need not be a physically correct observation. The archive also includes both an optional accelerated implementation and a Python path; agreement between these paths is not newly tested in this study.

\subsection{Historical presentation layer}
Three separately packaged HTML pages provide an analysis report, a project overview, and a test report. Their local JavaScript and JSON files are part of the presentation evidence. A page can have all its local assets present while lacking a record of the simulation and transformation steps that produced those assets. Accordingly, we audit each report bundle separately and do not merge its numerical statements with those of another page.

\section{Audit Method}
\label{sec:method}
\subsection{Study boundary and selection rule}
The unit of analysis is the local source archive supplied for manuscript preparation, examined on November 13, 2026. The archive has no local Git metadata. We therefore identify the examined material by an explicit selected-file inventory rather than assigning it an inferred upstream commit. Local directory names and timestamps are not used to infer original authorship or execution dates.

The selection rule includes Python and YAML files under the original simulation, root tooling, root tests, and workflow directories, excluding generated output and caches. It adds the simulation requirements, the root and two package README files, the reproducibility note, the source license, and every locally inventoried report asset. The new paper directory is outside these four selected source directories and is excluded from the source-archive inventory. The resulting inventory contains \SelectedFiles files. Each entry records a repository-relative path, a byte count, and a SHA-256 digest. The scope is selected-file integrity, not a full filesystem image.

Only the root generic tooling and its tests are copied into the paper's reproducibility bundle. Their bytes are retained unchanged, with the original license. No simulation module is imported or executed by the study scripts. A static AST traversal counts test-function definitions under the simulation tests; these counts are not pytest collection counts, since parametrization and runtime discovery can change the number of executable cases.

\subsection{File snapshots and inference limits}
For a selected file set $S$, let $b(p)$ denote the current bytes at relative path $p$. Define
\begin{equation}
L(S)=\operatorname{sort}_{p}\{(p,\lvert b(p)\rvert,\operatorname{SHA256}(b(p))):p\in S\}.
\label{eq:inventory}
\end{equation}
The study inventory records the digest of a compact, sorted-key JSON encoding of $L(S)$. The existing snapshot tool similarly records selected source files, specified inputs and artifacts, capture-time metadata, and a digest of the manifest payload before its identifier is added. This is an application-specific deterministic encoding; it is not claimed to implement a general JSON canonicalization standard.

Verification checks manifest structure, its digest, the selected source inventory, and recorded file bytes. Complete report verification additionally requires coverage of every asset discovered by the report inventory. A simplified condition is
\begin{equation}
V(M)=V_{\mathrm{structure}}\land V_{\mathrm{digest}}\land V_{\mathrm{selection}}\land V_{\mathrm{files}},
\qquad A_r\subseteq C_M,
\label{eq:verification}
\end{equation}
where $A_r$ is the inventoried local asset set for report $r$ and $C_M$ is the artifact-path set covered by its manifest. The report auditor also rejects a verified status when asset-resolution issues remain. Passing these checks establishes consistency with the supplied manifest within the stated scope. It does not establish the manifest's independent authenticity, the origin of a historical run, or the accuracy of reported results (Figure~\ref{fig:evidence}).

\begin{figure}[t]
  \centering
  \includegraphics[width=\linewidth]{figures/evidence_flow.pdf}
  \caption{Evidence available in this study and evidence not established by the archive. The contemporary tooling run is attributable to the saved environment and test records. It does not fill gaps in the historical execution chain.}
  \label{fig:evidence}
\end{figure}

\subsection{Static report inventory}
The report auditor starts from the three fixed HTML entry points. It follows explicit local runtime references in HTML and CSS and inventories packaged \code{assets/} and \code{\_shared/} directories, including files that may be loaded indirectly. It normalizes local paths, separates external or embedded URLs, and records missing files, paths outside the selected root, and unsupported resolution cases. Navigation hyperlinks are not treated as runtime dependencies.

The tool expects a report-specific \code{provenance.json} and assigns one of three labels:
\begin{description}[font=\normalfont,leftmargin=0pt]
  \item[\code{legacy\_unverified}:] the report-specific manifest is absent.
  \item[\code{integrity\_failed}:] a supplied manifest has inconsistent bytes, incomplete asset coverage, or unresolved inventory issues.
  \item[\code{snapshot\_verified}:] the manifest is valid and complete, and no inventory issues remain.
\end{description}
These labels describe the present snapshot check. They are not confidence scores or a ranking of physical-model quality.

The inventory does not execute JavaScript or fetch external resources. Runtime-generated references, remotely hosted libraries, and dynamically constructed CSS may therefore remain outside its observation boundary. ``No issues detected'' means no issues under this static rule, not proof that a page is fully portable or that every dependency has been discovered.

\subsection{Executed checks and saved records}
We run the copied root tooling suite with pytest and save its JUnit XML. These tests exercise metadata collection, snapshot validation, changed and missing files, selection boundaries, manifest handling, report-asset coverage, and command-line behavior. Passed, failed, and skipped cases are retained as separate outcomes, consistent with pytest's distinction between a skipped case and an executed success \citep{pytestskip}.

Ten additional scenarios use temporary files containing arbitrary text and a tiny HTML/JavaScript report. Each scenario records an explicit expected validity or status, the observed outcome, and diagnostic codes. The scenarios include unchanged files, five kinds of file or manifest mutation, a generated-output negative control, and three asset-coverage states. They characterize generic file-integrity behavior and exercise neither flight models nor performance metrics.

\section{Experimental Setup}
\label{sec:setup}
The tooling run uses CPython \AuditPython on Windows with pytest \AuditPytest. The environment record reads installed-distribution metadata without importing scientific packages. In that environment, pytest-xdist is absent, and the recorded pytest version is below the simulation requirements file's declared minimum. This does not prevent the independent root tooling suite from executing, but it prevents treating that environment record as proof that the simulator's dependency requirements were satisfied. We do not resolve or validate simulation-library compatibility.

The archived workflow specifies Windows/Linux and Python 3.11/3.12 jobs for the generic tooling. A workflow configuration describes intended checks, rather than proving that all jobs have completed. The present experiment is one local run, and no cross-platform or full-simulation result is inferred from that matrix.

Evidence is collected in four steps: static source and report inventory; execution of the copied generic test suite; controlled integrity probes; and generation of summary numbers and vector figures from the saved records. Appendix~\ref{app:commands} gives the commands. The manuscript's headline counts are generated into \code{evidence\_numbers.tex}; detailed audit and test records remain available alongside the manuscript. The original simulation files and historical reports are left unchanged.

\section{Results}
\label{sec:results}
\subsection{RQ1: implementation and test inventory}
Static inspection supports the division between simulation orchestration, module families, telemetry, and independent provenance tooling summarized in Section~\ref{sec:architecture}. The simulation test directory contains \SimulationTestFiles files defining \SimulationTestDefinitions test functions. These figures describe source inventory only. Those tests are not executed here, and we do not report their definitions as successful tests. Similarly, the presence of campaign and convergence utilities is an implementation observation rather than evidence of a reproduced campaign or convergence result.

\subsection{RQ2: report assets and provenance status}
Table~\ref{tab:reports} records the static report audit. The analysis, overview, and test report bundles contain 5, 7, and 1 locally inventoried files, respectively, for \ReportAssets distinct assets including the HTML entry points. None has the expected report-specific manifest. All three therefore remain \code{legacy\_unverified}. No missing-asset or other inventory issue is reported within the implemented static scope.

\begin{table}[t]
\caption{Contemporary static report audit. Asset counts include HTML entry points, packaged data, and local libraries. Zero detected issues does not establish historical execution provenance.}
\label{tab:reports}
\centering\small
\begin{tabular}{@{}lrrll@{}}
\toprule
Report & Assets & Issues & Manifest & Status \\
\midrule
Analysis & 5 & 0 & Absent & \code{legacy\_unverified} \\
Overview & 7 & 0 & Absent & \code{legacy\_unverified} \\
Test report & 1 & 0 & Absent & \code{legacy\_unverified} \\
\bottomrule
\end{tabular}
\end{table}

The distinction between availability and provenance is visible directly: all inventoried files can be present while none of the report bundles has a matching manifest. Running the report auditor in strict mode returns exit code 1. In this setting, the nonzero exit expresses unmet provenance requirements; it is not an observed simulation failure. No historical numeric values are promoted to newly verified measurements, and the pages are not assumed to describe a common run.

\subsection{RQ3: generic tooling behavior}
The saved JUnit record contains \TestsTotal collected cases: \TestsPassed passed, \TestsSkipped skipped, and \TestsFailed failed. Table~\ref{tab:tests} gives the suite-level breakdown. The skipped case concerns a packaged directory symlink escaping the root; directory symlink creation was unavailable in the local environment. Its intended behavior is therefore not established by this execution. Other boundary checks exercised by the passing cases remain narrower observations, not a substitute for that skipped check.

\begin{table}[t]
\caption{Executed independent tooling tests. These are root-level tooling cases, distinct from the unexecuted simulation test definitions.}
\label{tab:tests}
\centering\small
\begin{tabular}{@{}lrrr@{}}
\toprule
Suite & Passed & Skipped & Failed \\
\midrule
Environment inventory & 8 & 0 & 0 \\
Snapshot provenance & 28 & 0 & 0 \\
Static report audit & 12 & 1 & 0 \\
\midrule
Total & \TestsPassed & \TestsSkipped & \TestsFailed \\
\bottomrule
\end{tabular}
\end{table}

All \ProbeCount additional scenarios match their expected outcomes (Table~\ref{tab:probes}). A changed or deleted asset is rejected, a changed source file or an added selected source is detected, and editing manifest metadata without updating its digest invalidates the manifest. Adding a generated source-like file under an excluded output directory does not invalidate the selected-source snapshot. This negative control illustrates the selection boundary: a snapshot cannot detect changes to files that it deliberately excludes.

For the synthetic report, no manifest yields \code{legacy\_unverified}; a manifest covering only the HTML yields \code{integrity\_failed}; and a manifest covering both HTML and the discovered data asset yields \code{snapshot\_verified}. The last outcome concerns a newly constructed synthetic report. It does not change the status of any historical report in the archive. Figure~\ref{fig:results} summarizes the report and test observations.

\begin{table}[t]
\caption{Controlled integrity probes using nonphysical temporary files. All ten expectations were met. A ``valid'' snapshot remains subject to its selected-file scope.}
\label{tab:probes}
\small
\begin{tabularx}{\linewidth}{@{}Yp{0.38\linewidth}@{}}
\toprule
Scenario & Observed outcome \\
\midrule
Unchanged snapshot & Valid \\
Modified report data asset & Invalid: file changed \\
Deleted report data asset & Invalid: file unavailable \\
Modified selected source & Invalid: file changed \\
Added selected source & Invalid: source added \\
Edited manifest label without digest update & Invalid: manifest digest mismatch \\
Generated file in excluded output directory & Valid within selection \\
Absent synthetic report manifest & \code{legacy\_unverified} \\
Manifest omitting discovered data asset & \code{integrity\_failed} \\
Manifest covering all discovered report assets & \code{snapshot\_verified} \\
\bottomrule
\end{tabularx}
\end{table}

\begin{figure}[t]
  \centering
  \includegraphics[width=\linewidth]{figures/audit_results.pdf}
  \caption{Observed local artifact counts and independent test outcomes. The report asset counts measure inventory, while the test bars represent executed generic tooling cases. Neither panel measures simulation performance.}
  \label{fig:results}
\end{figure}
\FloatBarrier

\section{Discussion}
\subsection{What the archive can support}
The results support an architecture description and a contemporary artifact-audit report. They show that the selected files are inspectable, that the locally discovered report assets exist, and that the exercised generic tooling behaves as expected in the recorded environment. These statements are useful and independently narrower than a claim that the original simulation experiments have been reproduced.

A typical reporting error would combine the current source, an existing telemetry file, and a historical figure into one presumed run. This archive provides no verified association sufficient for that inference. The appropriate response is to retain the report as historical material with an explicit status, rather than infer missing execution metadata from a page title, file modification time, or current configuration.

The distinction also applies to test reporting. Test definitions, collected cases, successful executions, and a CI matrix are different observations. The \SimulationTestDefinitions simulation test definitions do not become passed cases because \TestsPassed separate tooling tests pass. Nor does a configured matrix become completed compatibility evidence. Keeping these observations separate makes the manuscript's conclusions easier to audit.

\subsection{A reporting contract for scientific software}
For a generic computational study, the evidence associated with a reported result should identify the selected source, configuration and inputs, execution environment, actual invocation, and resulting output artifacts. For stochastic computations, the recorded workflow should also identify the randomization state used by that computation. A later rendering step should retain its own inputs and output hashes. These are records of what happened, not a recipe for inferring what probably happened.

Such a record is stronger when captured during execution and retained with the artifacts. Post-hoc snapshots remain useful for detecting subsequent changes, but their capture environment must be labeled as the snapshot environment. They must not be represented as the environment that generated an older output. Likewise, a single manifest for a README cannot verify unrelated HTML reports, and a manifest covering an HTML entry point alone cannot verify unrecorded local data files.

This contract is a documentation recommendation, not an implemented historical recovery mechanism. The package does not backfill original runs, overwrite existing reports, or alter their provenance labels. No claim is made that the proposed record format is sufficient for independent physical validation or for an artifact-evaluation badge.

\section{Limitations and Threats to Validity}
The study examines one local archive and cannot estimate the frequency of these issues across scientific software. The selected-file rule omits generated directories and many non-source file types. Its \SelectedFiles-file inventory is therefore a declared subset, not a complete archive hash. The original files are not republished in the paper bundle, apart from the unchanged generic tooling and its license; verifying the original selected-file digests requires access to the examined archive.

The report inventory is static. It neither executes pages nor checks remotely hosted dependencies, and it may miss runtime-generated references. The ten controlled probes cover explicit scenarios, not all mutation patterns or adversarial behavior. The manifest digest is not a digital signature. A party able to replace both a manifest and its referenced files can produce a different self-consistent snapshot.

The executed tests are part of the archived tooling, rather than an independent implementation of its specification. Additional probes improve observability of selected cases but do not establish exhaustive correctness. One symlink case is skipped, only one local environment is exercised, and dependency metadata is not a compatibility test. The simulation itself, its original tests, numerical convergence, campaign statistics, physical fidelity, and historical output generation are not evaluated. The paper consequently offers no effectiveness, accuracy, or deployment claim for an interception system.

\section{Conclusion}
The \emph{Missile Intercept} archive supports a documented account of modular simulation software and a bounded audit of its current artifacts. Its three historical report bundles contain \ReportAssets inventoried local files but lack matching provenance manifests. A new run of the independent tooling yields \TestsPassed passed tests and \TestsSkipped skip, with all \ProbeCount controlled integrity probes matching expectation. These findings establish current software-evidence observations while leaving historical simulation claims unverified. Preserving that distinction is the principal result of the case study.

\section*{Artifact Availability and Authorship}
The accompanying paper folder includes LaTeX and BibTeX sources, compiled PDF, vector figures, saved audit records, and an unchanged copy of the generic audit tools and tests under their original MIT license. The original simulator archive is required only for repeating the static source/report inventory; the copied tooling tests and synthetic snapshot scenarios can be examined independently. The manuscript uses anonymous authorship during draft preparation. No author identity, affiliation, funding source, or original experimental history is inferred from the archive. No arXiv submission or external publication is performed by this study.

\bibliographystyle{unsrtnat}
\bibliography{references}

\appendix
\section{Reproduction Commands and Artifact Map}
\label{app:commands}
The paper directory is the working directory for the commands below. The first command expects the original archive as its parent; an alternative archive path may be supplied through \path{--source-root}. Evidence should be regenerated in a working copy to preserve the delivered records.

\begin{lstlisting}[basicstyle=\small\ttfamily,columns=fullflexible,keepspaces=true,breaklines=true,breakatwhitespace=true]
python scripts/prepare_evidence.py
cd reproducibility
python -m pytest -q tests -p no:cacheprovider --junitxml=../evidence/tooling_tests.xml
cd ..
python scripts/probe_tooling.py
python scripts/summarize_evidence.py
\end{lstlisting}
The pytest invocation above is one command. The probe script accepts \path{--source-root} for the optional strict audit of the original reports; when no archive is present, it retains the self-contained synthetic checks and records the unavailable strict-audit observation explicitly. Rendering the three figures additionally requires ReportLab. The manuscript can be compiled independently with standard pdfLaTeX and BibTeX, using the supplied figures and generated number file.

\begin{table}[h]
\caption{Claim-to-artifact map within the accompanying paper folder.}
\small
\begin{tabularx}{\linewidth}{@{}p{0.37\linewidth}Y@{}}
\toprule
Claim or record & File under the paper folder \\
\midrule
Selected original-file bytes & \path{evidence/source_inventory.json} \\
Simulation test-definition counts & \path{evidence/static_inventory.json} \\
Report assets and statuses & \path{evidence/report_audit.json} \\
Capture-time environment & \path{evidence/environment.json} \\
Executed tooling cases and skip reason & \path{evidence/tooling_tests.xml} \\
Ten controlled probe observations & \path{evidence/integrity_probes.json} \\
Aggregated counts and manuscript macros & \path{evidence/summary.json}, \path{evidence_numbers.tex} \\
Unchanged tooling and tests & \path{reproducibility/tools/}, \path{reproducibility/tests/} \\
Detailed evidence interpretation & \path{evidence/CLAIM_MAP.md} \\
Delivered package integrity & \path{package_manifest.json} \\
\bottomrule
\end{tabularx}
\end{table}
\end{document}
